\chapter{相关理论基础}

\section{电机故障类型与机理}

电机在运行过程中可能发生多种类型的故障，根据故障部位和原因的不同，主要分为以下几类：

\subsection{轴承故障}

轴承是电机中最容易发生故障的部件之一。常见的轴承故障包括外圈故障、内圈故障、滚动体故障和保持架故障。轴承故障的振动信号具有周期性冲击特征，其故障特征频率可由以下公式计算：

\begin{equation}
	f_{BPFI} = \frac{N_b}{2} f_r \left(1 + \frac{d}{D} \cos\alpha\right)
	\label{eq:bpfi}
\end{equation}

其中，$N_b$为滚动体数量，$f_r$为转频，$d$为滚动体直径，$D$为轴承节径，$\alpha$为接触角。

\subsection{转子故障}

转子故障主要包括转子不平衡、转子不对中和转子弯曲等。转子不平衡是最常见的转子故障，其振动特征表现为1倍转频分量幅值增大。转子不对中则会产生2倍转频分量，并可能激发高次谐波。

\section{深度学习基础理论}

\subsection{卷积神经网络}

卷积神经网络（CNN）是深度学习中最重要的网络结构之一，特别适合处理具有网格结构的数据。一维CNN的基本操作包括：

（1）卷积运算：第$l$层第$i$个特征图的第$j$个元素计算为

\begin{equation}
	x_{i,j}^{(l)} = f\left(\sum_{m=1}^{M} \sum_{k=1}^{K} w_{i,m,k}^{(l)} x_{m,j+k-1}^{(l-1)} + b_i^{(l)}\right)
	\label{eq:conv1d}
\end{equation}

其中，$w_{i,m,k}^{(l)}$为卷积核权重，$b_i^{(l)}$为偏置，$f(\cdot)$为激活函数，$K$为卷积核大小。

（2）池化运算：最大池化和平均池化是两种常用的池化方式，用于降低特征图维度，增强特征的平移不变性。

（3）全连接层：将卷积和池化层提取的特征进行整合，输出分类结果。

\subsection{迁移学习}

迁移学习是利用源域$\mathcal{D}_S$上学习到的知识来帮助目标域$\mathcal{D}_T$上的学习任务。当$\mathcal{D}_S \neq \mathcal{D}_T$或$\mathcal{T}_S \neq \mathcal{T}_T$时，迁移学习能够有效解决目标域数据不足的问题。

基于预训练微调的迁移学习策略包括两个阶段：（1）预训练阶段：在源域数据上训练模型，学习通用的特征表示；（2）微调阶段：冻结部分网络层，在目标域数据上微调剩余层，适应目标任务。

\section{本章小结}

本章介绍了电机故障诊断和深度学习的相关理论基础，包括电机常见故障类型及其振动特征、卷积神经网络的基本结构和迁移学习的基本原理，为后续章节的方法研究奠定了理论基础。